\section{A curve through the prescribed coefficients}\label{sec:descent}

An exact circuit gives a single choice of parse-table entries that realizes
all the coefficients of $P_n^*$. A border circuit gives a point in the
closure of the coefficient image. We realize that point as the value of the
whole table along one branch of an algebraic curve. At the chosen place,
the table entries may have poles, but every output coefficient is regular.
For example, the map $(x,y)\mapsto(x,xy)$ takes $(t,t^{-1})$ to
$(t,1)$. The value $(0,1)$ at $t=0$ has no finite preimage.
For the parse table, the processors in Section~\ref{sec:cartier} will
need equations for the curve and a short procedure for evaluating them.

Fix the compressor identities of Section~\ref{sec:compr} in the given
characteristic. Suppose that $\bsz_F^{[n]}(P_n^*)\le m_0$, with $m_0,M,R$
as in Section~\ref{sec:parse}. In this and the next two sections,
\emph{polynomial} means $M^{O(1)}$ and \emph{single exponential} means
$2^{M^{O(1)}}$. The exponents may depend on the fixed identities and are
uniform in $n$. Recall that $M^C=o(R)$ for every constant $C$.

Let $c=(c_\nu)_{|\nu|=n}$ be the integer coefficient vector of $P_n^*$
and write the table of Lemma~\ref{lem:parse-tables} at size $m_0$ as
\[
 T_X(Z)=\sum_{|\nu|=n}Y_\nu(X)Z^\nu,\qquad X=(X_1,\ldots,X_h).
\]
Here $Z$ lists all $2n$ paired input variables and $\nu$ ranges over
\emph{every} monomial of degree $n$. The closure of the coefficient
image contains $c$, interpreted over the algebraic closure of the prime
field as in Section~\ref{sec:parse}. The number $h\ge1$ is polynomial,
and each $Y_\nu$ has degree at most $d=2n-1$ and, when nonzero,
coefficient $\ell^1$-norm of polynomial bit length. The clock defining
$U_n^*$ gives the same bit bound for the entries of $c$.

\subsection{Specialization in characteristic zero}

To read a boundary value by residues we will work in positive
characteristic. In characteristic zero we can choose a prime separately
for each $n$. The prime must preserve the closure membership and have
bit length small enough for the processors.

\begin{lemma}\label{lem:descent-prime}
If $c\in\overline{Y(\mathbb A^h)}$ over $\overline{\mathbb Q}$, there is a
prime $p$ of polynomial bit length for which the same inclusion holds over
$\overline{\mathbb F}_p$. The fixed tensor identities can be specialized
at the same prime, with coefficients of bounded extension degree.
\end{lemma}

We use the following degree estimates over algebraically closed fields
(see \cite[Ch.~I, \S7]{Hartshorne} and
\cite[Lemma~2 and Proposition~3]{Heintz}).
For an affine variety we take the degree of its projective closure.
A linear image of an irreducible variety has closure of degree at most
the degree of the variety. For a rational map on an open subset of
$\mathbb A^h$ whose numerators and common denominator have degree at most
$e\ge1$, the degree of the image closure is at most $e^h$. Indeed, a
general linear section of complementary dimension meets the image closure
in $D'$ distinct points of the image, where $D'$ is the degree of the
closure. In the domain its pullback has at least one component over
each point. B\'ezout bounds the number
of components by $e^h$, even in mixed dimensions. At each cut of degree
at most $e$, the sum of $(\deg V)e^{\dim V}$ over the components $V$
cannot increase. In particular, linear sections never increase the sum
of component degrees.

We will also use that a closed irreducible set of degree $D$ can be
separated from any point outside it by an equation of degree at most $D$.
For a set of dimension $b$ and codimension at least two, project its
projective closure to $\mathbb P^{b+1}$ from a linear center disjoint from
its join with the point. The image misses the projected point and has
dimension $b$, since a positive-dimensional projective fiber would meet the
center. Its hypersurface equation pulls back to the required separator.
In codimension one the defining equation suffices. For a reducible set,
multiply the separators for its components.

\begin{proof}[Proof of Lemma~\ref{lem:descent-prime}]
The image closure has degree at most $D=d^h$ in every characteristic.
It therefore suffices to preserve the vanishing at $c$ of all relations
of degree at most $D$. Record substitution $y\mapsto Y(X)$ on these
polynomials by an integer matrix $A$. Its column for $y^\eta$ is the
coefficient vector of $Y(X)^\eta$ in $X$-monomials of degree at most
$dD$. The row $u=(c^\eta)_{|\eta|\le D}$ records evaluation at $c$.
The rational kernel of $A$ consists of the relations, so the closure
hypothesis puts $u$ in the rational row span.

There are only $\binom{h+dD}{h}$ rows, a single exponential, regardless
of the number of columns. Submultiplicativity of the coefficient norm
bounds the bit length of each matrix entry by a single exponential.
Choose a row basis and a square minor on those rows with nonzero
determinant $\Delta$. Cramer's rule expresses $u$ in the basis with
denominators dividing $\Delta$. The expression obtained on the minor's
columns holds on all columns because $u$ is in the row span. Hadamard's
inequality gives a singly exponential bound on the bit length of $\Delta$.

The constants of the fixed identities lie in a fixed number field and
are integral after inverting a fixed integer. Outside a finite list of
primes we may reduce at any prime ideal over $p$, preserving both the
identities and any required nonvanishings. The residue degree is at most
the number-field degree. Some prime $p$ of polynomial bit length lies
outside this list and does not divide $\Delta$, since the logarithm of
the product of all primes up to $x$ is asymptotic to $x$. Reducing the row
expression now puts $u\bmod p$ in the row span of $A\bmod p$. It therefore
annihilates every relation of degree at most $D$ even over
$\overline{\mathbb F}_p$. The separator bound proves the inclusion.
\end{proof}

In the diagonal argument the compressor identities were chosen over the
prime field in Section~\ref{sec:compr}. Their characteristic-zero
constants are rational, so the specialization above places them in
$\mathbb F_p$. The original layouts therefore give the same contraction
identities over $\mathbb F_p$, whether or not this new characteristic
admits ports. In fixed positive characteristic we keep the original
prime and constants. In both cases write $\mathcal K=\overline{\mathbb F}_p$
and let $E_0$ be a finite field containing the constants, of degree over
$\mathbb F_p$ bounded independently of $n$. For the prime-field choices
we can take $E_0=\mathbb F_p$.

\subsection{Equations for the branch}

Lehmkuhl and Lickteig for tensors, and B\"urgisser for circuits, used
the degree of the graph of a parameter map to bound the order of
approximation \cite{LehmkuhlLickteig,BurgisserCorrection}. Guo, Saxena and
Sinhababu give the corresponding classical bounded-degree approximation
for approximate polynomial satisfiability over arbitrary fields
\cite[Theorem~4.2]{GSS}. Here we need a separable curve parameter $t$
and equations that the processors can evaluate in polynomial time.

\begin{proposition}\label{prop:descent-curve}
There are a finite extension $E$ of $E_0$, with $\log_2|E|$ polynomial, and
$h$ polynomials $g_i(X,t)\in E[X,t]$ of polynomial degree with the following
properties. Their common zero set over $\mathcal K$ has an irreducible curve
component of singly exponential degree, dominant over the $t$-line, on which
$J=\det(\partial g_i/\partial X_j)$ is generically nonzero. The function
field of this curve has a place above $t=0$ at which all $Y_\nu$ are regular
with residue $c_\nu$. With polynomial advice, each $g_i$ can be evaluated
in polynomial bit time over $E$ or a specified extension of polynomial bit
length, given its field representation and the embedding of $E$.
\end{proposition}

\begin{proof}
Close the graph of $Y$ in $\mathbb P^h\times\mathbb A^{\{\nu:|\nu|=n\}}$
using homogeneous parameter coordinates $(1:X)$. The projection to output
space is closed, so the fiber above $c$ is a nonempty projective set,
possibly entirely at parameter infinity. On each standard parameter
chart the graph closure has degree at most $(d+1)^h$. Indeed, its dense
finite-parameter part is parametrized by $X$ with chart coordinates
given by ratios of affine linear forms and output coordinates $Y(X)$.
With the chart denominator common to all coordinates, the numerators
have degree at most $d+1$. The fiber above $c$ is a linear section on
each chart. Counting each projective component in a chart meeting it
bounds the sum of the fiber component degrees by $(h+1)(d+1)^h$.

Cut the projective fiber by parameter hyperplanes until it is finite.
At each cut avoid the hyperplanes containing a positive-dimensional
component. Every such component meets the chosen hyperplane by
projectivity, so the last intersection is still nonempty. Each component
excludes at most a proper linear subspace of coefficients, and the
number of exclusions at a cut is bounded by the degree sum above.
Fix a point $P$ of the finite intersection and a standard chart containing
it. Denote the graph closure in this chart by $G$, write $(z,y)$ for the
chart and output coordinates, and let $\ell_1(z),\ldots,\ell_b(z)$ be
the hyperplanes used in the cuts. On $G$, the common zeros of the forms
$y_\nu-c_\nu$ and $\ell_j$ now make up a finite set containing $P$.

We need $h$ equations for a finite intersection, with short descriptions
even when the table has many output coordinates. Encode each $\nu$ by
its base-$(n+1)$ digits. With $B=(n+1)^{2n}$, put
\[
 e(\nu)=\sum_{j=1}^{2n}\nu_j(n+1)^{j-1},\qquad
 D_a(z,y)=\sum_\nu a^{e(\nu)}(y_\nu-c_\nu)
             +\sum_{j=1}^b a^{B+j}\ell_j(z).
\]
On any surviving positive-dimensional component, at least one original
form is nonzero at some point. At that point $D_a$ is a nonzero
polynomial in $a$ of degree at most $B+b$, because all the exponents
are distinct. Avoiding its roots for each such component cuts every
positive dimension by one. These are linear cuts on $G$, so at each
step there are at most $\deg G$ components. After $h$ choices
$a_1,\ldots,a_h$, the common zeros of $D_i=D_{a_i}$ form a finite
set containing $P$.

We now deform this intersection. For a vector $v$ to be chosen, impose
\begin{equation}\label{eq:descent-deform}
                    D_i(z,y)+t(z_i-v_i)=0\qquad(1\le i\le h).
\end{equation}
At $t=0$ these equations retain $P$. Over $\overline{\mathcal K(t)}$,
they cut out the fiber over $v$ of the affine linear map
$\psi_t(z,y)=z+t^{-1}D(z,y)$, where $D=(D_1,\ldots,D_h)$.
Let $U\subset G$ be the finite-parameter part of this chart. It is an
open subset of $z$-space, with $y=Y(X(z))$. Write
$\widetilde D(z)=D(z,Y(X(z)))$. The determinant
$\det(tI+\partial\widetilde D/\partial z)$ is monic in $t$ over
$\mathcal K(G)$. Thus its vanishing on $U$ cuts out a set of dimension
less than $h$ over the generic $t$, even if the derivatives of all
outputs vanish in characteristic $p$.

Choose $v$ outside the closures of the images under $\psi_t$ of this
critical locus and of $G\setminus U$. They lie in a hypersurface of
singly exponential degree. Indeed, $G\setminus U$ is empty or cut out
by one chart coordinate. Each $\widetilde D_i$ has numerator of degree
at most $d+1$ and denominator a common linear form to power $d$.
Clearing the denominators in the determinant gives a polynomial of
degree $O(h(d+1))$ in $z$. By B\'ezout the sum of the degrees of the
two exceptional loci, after closure, is $O(h(d+1))\deg G$. Taking
linear images and multiplying their component separators gives a
hypersurface of at most this degree. Consequently every point
of the generic fiber has finite parameters and invertible $z$-Jacobian.

All hyperplane coefficients, the $a_i$ and $v$ can be chosen in one
extension $E$ of $E_0$ whose elements have polynomial bit length.
At every stage the product of the equations to be avoided has singly
exponential degree, uniformly in the earlier choices. A grid in a field
of size larger than that degree contains a point where the product is
nonzero, including for the exceptional hypersurface over
$\overline{\mathcal K(t)}$. Since $\log_2|E_0|$ is polynomial, we
can take a finite extension of this size with $\log_2|E|$ polynomial.

In $G\times\mathbb A^1$ the $h$ equations in
\eqref{eq:descent-deform} have local dimension at least one at $(P,0)$
by Krull's dimension theorem. Every component through that point of positive
dimension dominates the $t$-line, since the special fiber is finite.
The generic fiber has dimension at most zero by the invertible Jacobian,
so each such component is a curve. Choose one and take a point above
$(P,0)$ on its normalization. The functions $y_\nu$ are regular at
this place, with values $c_\nu$. On the dense finite-parameter part
they equal $Y_\nu(X)$, which therefore have the same values even if
some $X_j$ have poles.

It remains to give the equations in the original table entries $X$.
Write the chart coordinates as $z_i=u_i(X)/\delta(X)$ with $u_i$ and
$\delta$ affine linear. Clearing the denominator gives
\[
 g_i(X,t)=\delta(X) D_i\bigl(z,Y(X)\bigr)
                  +t\bigl(u_i(X)-v_i\delta(X)\bigr).
\]
This is polynomial because $D_i$ is affine linear in $z,y$. On the
finite-parameter chart the change of coordinates from $z$ to $X$ has
invertible Jacobian. At a common zero multiplying the equations by
$\delta$ scales each Jacobian row by a unit. Hence the affine closure
of the chosen curve in $X,t$ is a component with $J\ne0$ generically.
Its degree is at most $(d+2)^h$ by B\'ezout, since the equations have
degree at most $d+2$. The relations $dg_i=0$ express every $dX_j$ as
a multiple of $dt$. The curve has one-dimensional function-field
differentials over the perfect field $\mathcal K$, so $dt$ spans them.

For the evaluator, the large sum of output coordinates in each $D_i$
is just one table evaluation. For $a=a_i$,
\[
 \sum_\nu a^{e(\nu)}Y_\nu(X)
     =T_X\bigl(a^{(n+1)^0},\ldots,a^{(n+1)^{2n-1}}\bigr).
\]
Store the corresponding offset $\sum_\nu a^{e(\nu)}c_\nu$ as one
element of $E$. The $a_i$, hyperplane coefficients, $v$, offsets and
field arithmetic data take polynomially many advice bits. The exponents
have polynomial bit length, and Lemma~\ref{lem:parse-tables} supplies
the polynomial-size table recursion with weights $X$. Binary powering,
that recursion and the affine forms $u_i,\delta$ thus evaluate $g_i$
in polynomial bit time, using the cleared expression even at
$\delta(X)=0$. This evaluator uses only field data for the equations.
The geometric point $P$ and its branch may be defined over larger fields.
\end{proof}
